\section{Stakeholders}

\begin{frame}{Stakeholders}
\mytext{
    Someone who has a stake in something
    
    \medskip
    
    $\rightarrow$ who has something to win (or loose)

   \medskip
   
    \begin{itemize}
        \item They who want it: $\rightarrow$ \textit{initiators}
        \item They who do the work: $\rightarrow$ \textit{executors}
        \item They who will use it: $\rightarrow$ \textit{users} 
        \item They who are affected in other ways: $\rightarrow$ \textit{others}
    \end{itemize}
}
\end{frame}

\begin{frame}{Who wants it?}
    \mytext{
    \begin{itemize}
        \item Authorities
            \begin{itemize}
                \item different levels of Government
                \item different ministries
            \end{itemize}
        \item Business people
            \begin{itemize}
                \item those who finance the project
                \item those who realise the project
                \item those who operate the result
            \end{itemize}
        \item Users
            \begin{itemize}
                \item corporations
                \item groups
                \item individuals
            \end{itemize}
        \item Affected non-users
            \begin{itemize}
                \item may oppose it
                \item may be organised
            \end{itemize}
    \end{itemize}
    }
\end{frame}

\begin{frame}{Who does it?}
    \mytext{
    \begin{itemize}
        \item The project manager
        \item The project management office
        \item Contractors
        \item Sub-contractors
        \item Suppliers
        \item Workers and employees
        \item The steering committee
        \item The project owners
        \item Public services
        \item Money giver? Banks?
    \end{itemize}
    }
\end{frame}

\begin{frame}{Who will use it?}
    \mytext{
    \begin{itemize}
        \item Corporations
        \item Groups
        \item Individuals
        \item Providers
        \item Customers
        \item Machines, vehicles
        \item Robots
        \item Children
    \end{itemize}
    }
\end{frame}

\begin{frame}{Who else are affected?}
    \mytext{
    \begin{itemize}
        \item Regulators, authorities
        \item Neighbours
        \item Competitors
        \item Future generations
        \item Nature
        \item Markets
        \item Gangsters
        \item Politicians
        \item The Internet
        \item ...
    \end{itemize}
    }
\end{frame}


\begin{frame}{Stakeholder Journey}
    \mytext{
    \begin{itemize}
        \item How do they get involved?
        \item First impression?
        \item First reaction?
        \item Positive side?
        \item Negative side?
        \item What is their stake?
        \item Where does the path split?
        \begin{itemize}
            \item when are they convinced?
            \item when do they change their stand?
            \item when do they `up their ante'?
            \item when do they loose interest?
        \end{itemize}
        \item What narrative do you tell them?
        \item What narrative do you want them to tell?
    \end{itemize}
    }
\end{frame}

% \begingroup
%     \setbeamercolor{background canvas}{bg=violet!20}
% \begin{frame}{Task}
%     Describe one (ideal or typical) stakeholder of your case study!
%     \begin{itemize}
%         \item Tell the story
%         \item What was / is in for them?
%     \end{itemize}
% \end{frame}
% \endgroup

\begin{frame}{Stakeholder Matrix}
    \mytable{
    
    \renewcommand{\arraystretch}{2}
    \begin{tabular}{|r|c|c|}
    \hline
         & internal & external \\
         \hline
         initiators & & \\
         \hline
         executors & & \\
         \hline
         affected & & \\
         \hline
    \end{tabular}
    }
\end{frame}

\begin{frame}{Stakeholder Matrix}
        \showmytikz{../shared/tikz/forest-stakeholders-Suchanek-core}
\end{frame}
    
\begin{frame}{Stakeholder Matrix}
        \showmytikz{../shared/tikz/forest-stakeholders-suchanek1}
\end{frame}

\begin{frame}{Stakeholder Matrix}
    
        \showmytikz{../shared/tikz/forest-stakeholders-suchanek2}
\end{frame}

\begin{frame}{Stakeholder Matrix}
        \showmytikz{../shared/tikz/forest-stakeholders-suchanek-general}
\end{frame}

\begin{frame}{More on stakeholders}
\mytext{
    \begin{itemize}
        \item directly vs. indirectly involved stakeholders
        \item owners / invested stakeholders
        \item interests of people involved
        \item roles people play
        \item \textbf{actual persons}
        \item different perspectives
        \item different expertise
        \item different power
        \item different personal position
    \end{itemize}
}
\end{frame}


\begin{frame}{What is in it for me?}
\mytext{
    People usually want to gain...
    \begin{itemize}
        \item money?
        \begin{itemize}
            \item money quick?
            \item money long term?
        \end{itemize}
        \item fame?
        \item security?
        \item idealistic win?
        \item something for those near to them?
    \end{itemize}
    
    \medskip
    Are they rational? Does that matter?
}
\end{frame}

\begin{frame}{What do I put into it?}
\mytext{
    \begin{itemize}
        \item money?
        \begin{itemize}
            \item when, to whom, and how
            \item to which conditions
        \end{itemize}
        \item publicity?
        \item expertise?
        \item connections?
        \item work hours 
        \begin{itemize}
            \item my own?
            \item my team members?
        \end{itemize}
    \end{itemize}
}
\end{frame}

\begin{frame}{Who can stop it?}
\mytext{
    \begin{itemize}
        \item Ask: who can do it?
        \item but also: who can spoil it?
        \item Questions of power and influence
    \end{itemize}
}
\end{frame}


\begin{frame}{What can I do to stop it?}
\mytext{
    \begin{itemize}
        \item withdraw my support?
        \item talk bad about it / you?
        \item not do what I am supposed to do?
        \item confuse people?
        \item involve more people?
        \item go to court?
        \item ...
    \end{itemize}
}
\end{frame}

 
\begin{frame}{Stakeholder Map / Team Map}  % depends on: forest

    \begin{forest}
    for tree={
        font=\footnotesize,
        grow = east,
        parent anchor = east,
        child anchor = west,
        forked edge,
        inner sep=2pt,
        anchor=base west
            }
    [\parbox{20ex}{\centering Project: \\ Acquisition of another company}
        [3rd parties [...][...]]
        [Account managers[...][...]]
        [Human resources[...][...]]
        [The media[...][...]]
        [Business divisions[...][...]]
        [Marketing[...][...]]
        [Technology[...][...]]
        [Finance[...][...]]
    ]
    \end{forest}
    
\end{frame}


\begin{frame}{Stakeholder Matrix}
    \showmytikz{../shared/tikz/tikz-stakeholdermatrix1}
    \mycomment{
    
    \par
    Source: Roberts p. 40
    
    }
\end{frame}

\begin{frame}{Stakeholder Matrix}
    \showmytikz{../shared/tikz/tikz-stakeholdermatrix2}
    \mycomment{
    
    \par
    Source: Roberts p. 40
    }
\end{frame}



\begin{frame}{No Man is an Island!}
\mytext{
    \begin{itemize}
        \item Every project has multiple stakeholders
        \item direct stake holders
        \item owners
        \item steering groups
        \item project management team
        \item team members
        \item supervisors
        \item consultants
        \item potential clients
        \item 3rd parties, indirect stakeholders
        \item etc.
    \end{itemize}
}
    
\end{frame}


% \begin{frame}[allowframebreaks]{Environment}
%     \begin{tikzpicture}
%         \draw (0,0) rectangle (5,4);
%         \node[draw] at (2.5,2) {Project};
%     \end{tikzpicture}
    
%     \framebreak
    
%     \begin{tikzpicture}
%         \draw (0,0) rectangle (5,4);
%         \draw (1,1) rectangle (4,3);
%         \node[draw] at (2.5,2) {Project};
%     \end{tikzpicture}
    
%     \framebreak
    
%     \begin{quote}
%         A successful project is best characterised by the change it delivers. \cite[p.12]{Roberts}
%     \end{quote}
% \end{frame}


% \begin{frame}{Product Flow Chart (part)}
%     \tikzstyle{decision} = [diamond, draw, fill=blue!20, 
%         text width=4.5em, text badly centered, node distance=2cm, inner sep=0pt]
%     \tikzstyle{block} = [rectangle, draw, fill=blue!20, 
%         text width=5em, text centered, rounded corners, minimum height=4em]
%     \tikzstyle{line} = [draw, -latex']
%     \tikzstyle{cloud} = [draw, ellipse,fill=red!20, node distance=3cm,
%         minimum height=2em]
        
%     \begin{tikzpicture}[scale=0.2, node distance = 2.2cm, auto]
%     \footnotesize
%         % Place nodes
%         \node [block] (case) {business case};
%         \node [block, below of=case] (gov) {governance document};
%         \node [block, below left of=gov] (use) {user requirements document};
%         \node [block, below right of=gov] (sol) {solution design document};
%         \node [block, below right of=use] (sel) {selected factory};
%         % Draw edges
%         \path [line] (case) -- (gov);
%         \path [line] (gov) -- (use);
%         \path [line] (gov) -- (sol);
%         \path [line] (use) -- (sel);
%         \path [line] (sol) -- (sel);
%     \end{tikzpicture}
    
%     Source: \cite[p. 169]{Roberts}
    
% \end{frame}

% \begin{frame}{Hierarchy}
%     \smartdiagramset{back arrow disabled=true}
%     \smartdiagram[flow diagram]{Project Sponsor, Project Manager, Project Team}

%     \vfill
%     The reality is more complicated!
% \end{frame}




% \begin{frame}[allowframebreaks]{Project Live Cycle}
%     \begin{adjustbox}{width=\textwidth}
%         \smartdiagram[sequence diagram]{
%         Initiali\-sation,
%         Planning,
%         Execution,
%         Closure
%         }
%     \end{adjustbox}
% \end{frame}


\begin{frame}{Stake holder number over time}
    \mytext{
    \begin{itemize}
        \item Number grows tendentially over times
        \item `From top to bottom'
        \item Some will quit having a stake after a while (payoff)
        \item After the end of project...?
        \item `And then the lawyers get involved' 
    \end{itemize}
    }
\end{frame}

\begin{frame}{Stakeholder Involvement Strategy}
\mytable{
    \scriptsize
    \begin{tabular}{>{\RaggedRight}p{18ex}>{\RaggedRight}p{33ex}>{\RaggedRight}p{20ex}}
    \toprule
     Stakeholder \& Role & Involvement & Strategy \\
     \midrule
     John Doe, Owner & 
     Came up with the project. Millionaire from the US. Has already found funding. Philanthropic motivation. & Report monthly. Provide figures.\tabularnewline
     \addlinespace
     Jane Doe, Manager & Wife of John Doe, in charge of making the project happen. Visual artist. &  Discuss PR and visuals. Give ceremonial tasks. \\
     \addlinespace
     Rick Mayer,\\ Personal Assistant to the Does & The man in charge of everything. Trained in business administration. Not qualified in content matters of the project.  & Speak to daily or as required. Provide detail on progress and requirements. \\
      ... & ... & ... \\
      ... & ... & ... \\
     \bottomrule
    \end{tabular}
}
\end{frame}


% \input{break}

\begin{frame}{From Project to Product: `Port'}
    \showmytikz[caption={From Project to Product: `Port' (Stakeholders)},label=fig:port-stakeholders]{../shared/tikz/tikz-port.tex}
\end{frame}


% \begin{frame}[allowframebreaks]{Project Kickoff}
%     First kickoff... then kickoff.. then... wait?
%     \medskip
%     What exactly is a `kickoff'?
%     \medskip
%     \begin{itemize}
%         \item Rewarding of contract
%         \item Celebrating closure of contract 
%         \item Public declaration
%         \item Ground breaking ceremony
%         \item Topping out ceremony
%     \end{itemize}
    
%     \framebreak
    
    
    
    
%     \begin{itemize}
%         \item define clear milestones
%         \item define clear outcomes of meetings
%         \item make decisions public
%     \end{itemize}
% \end{frame}


% \begin{frame}[allowframebreaks]{Coming up with a Project}  % depends on: smartdiagram, adjustbox
%     \begin{adjustbox}{width=\textwidth}
%         \smartdiagramset{back arrow disabled=true}
%         \smartdiagram[flow diagram:horizontal]{Notice, Understand, Find, Formulate, Convince}
%     \end{adjustbox}
%     \begin{description}
%         \item[Notice]  {a problem / an opportunity},
%         \item[Understand] {causes / the situation}, 
%         \item[Find] {a solution / a business case}, 
%         \item[Formulate] {the business / project plan},
%         \item[Convince] {stakeholders},
%         \item[Plan] {tasks, budget, team, time},
%         \item[Do] {work, controlling, reporting, correcting}
%         \item[Deliver] {product, result, review}
%     \end{description}
    
%     \framebreak
    
%     \begin{adjustbox}{width=\textwidth}
%         \smartdiagramset{back arrow disabled=true}
%         \smartdiagram[flow diagram:horizontal]{receive notice, analyse, formulate, begin}
%     \end{adjustbox}
    

%     \begin{description}
%         \item[Receive notice]  {from project sponsor},
%         \item[Understand] {stakeholders, their role, motivation etc.}, 
%         \item[Find] {a strategy}, 
%         \item[Formulate] {the project plan},
%         \item[Convince] {stakeholders},
%         \item[Do] {planning, work, controlling, reporting, correcting}
%         \item[Deliver] {product, result, review}
%     \end{description}

% \end{frame}




% \begin{frame}[allowframebreaks]{Typical Main Project Phases (Suchanek)}
%     \resizebox{!}{0.8\textheight}{
%         \input{tikz/tikz-chain-projectphases-details}
%     }

%     \framebreak
    
%     \resizebox{\textwidth}{!}{
%     \begin{meinganttq}{5}
%         \inserttable{gantts/mains.csv}
%         \showmilestones{gantts/mains.csv}
%     \end{meinganttq}
%     }
% \end{frame}

% \begin{frame}{Project Launch per Kerzner}

% \resizebox{!}{0.4\textheight}{
% \parbox{\textwidth}{
%     In the launch phase, customer (= project sponsor), project manager (pm), subject matter experts (sme's) and functional leads (leads) are involved. 
%     \begin{description}
%     \item[Phase 1] - Initialisation
%     \begin{itemize}
%         \item Kickoff-Meeting with customer
%     \end{itemize}
%     \item[Phase 2] - Preliminary planning
%     \begin{itemize}
%         \item PM, SMEs, Leads: Scope Management plan
%         \item Perliminary Scope Statement
%     \end{itemize}
%     \item[Phase 3] - Detailed planning
%     \begin{itemize}
%         \item PM, SMEs, Leads: \textbf{Work Breakdown Structure (WBS)} development
%         \item WBS dictionary
%     \end{itemize}
%     \item[Phase 4] Scope baseline
%     \begin{itemize}
%         \item Finalisation of governing documents\\ ("Lastenheft", product requirement documentation)
%     \end{itemize}
%     \end{description}
    
%     Source: \cite[p. 158]{KerznerHarold2017Pm:a}
% }
% }    
% \end{frame}

% \begin{frame}{Project Kick off}
% Example: project `play football'
%     \begin{itemize}
%         \item Before you kick the ball...  
%         \item Players must already be committed
%         \item Rules and commitment must be clear
%         \item You don't want an embarrassment
%         \item Your ideas must already be their ideas
%         \item Stakeholders are a team
%         \item Ceremonial event, team building
%         \item Believe in success!
%     \end{itemize}
% \end{frame}



% \begin{frame}{The Objectives of a Project}

%     Well defined objectives are prerequisite. BUT
%     \begin{itemize}
%         \item objectives may be subject to change.
%         \item different stake holders may have different views.
%     \end{itemize}
    
    
%     Defined roles and institutions are much less easy to change.
    
% \end{frame}


% \begin{frame}{Project Governance Document (PGD)}
%     \begin{itemize}
%         \item A PGD is also the `constitution' of the project
%         \item you must identify / find stakeholders
%         \item you have to create institutions
%         \item define obligations / duties
%     \end{itemize}
% \end{frame}

% \begin{frame}{From Knowledge to Action}

%     \begin{quote}
%         \flushright Der Mensch kann viel, wenn er nur muss\\
%         \flushright---unknown
%     \end{quote}

%     \begin{itemize}
%         \item to have the right to
%         \item to know how to
%         \item to have reason to
%         \item to want to
%         \item to have what it takes to
%         \item to have the duty to
%     \end{itemize}
% \end{frame}

% \begin{frame}{Obligations / Duties}
%     \begin{itemize}
%         \item bind yourself to your word
%         \item bind yourself to your supporters
%         \item bind yourself to your goals
%         \item set yourself s.m.a.r.t. goals
%         \item promise stakeholders s.m.a.r.t. goals
%         \item set team members and contractors s.m.a.r.t. goals
%         \item report and control
%     \end{itemize}
% \end{frame}


% \begin{frame}{Idea, initialisation}
%     \begin{itemize}
%         \item Why do it?
%         \item What to do to achieve what?
%         \item Who wants it?
%         \item Who pays for it?
%         \item Who might oppose it?
%         \item What will we need?
%         \item (What are the opportunities?)
%         \item (What are the risks?)
%         \item (What are the drawbacks?)
%         \item Whom do we need to convince?
%         \item To whom do we have to communicate?
%     \end{itemize}
% \end{frame}

% \begin{frame}{What Counts}
%     \begin{itemize}
%         \item Relevant for the Project Owner / Sponsor:
%         \begin{itemize}
%             \item opportunities and risks of the project
%             \item business case, business strategy
%             \item \textbf{why}
%         \end{itemize}
%         \item For the Project Manager:
%         \begin{itemize}
%             \item \textbf{who} wants this? 
%             \item \textbf{who} does it?
%             \item \textbf{what do they actually want?}
%             \item \textbf{with what }can we do it?
%             \item opportunities and risks \textbf{of tasks / sub-projects}
%             \item stakeholder satisfaction
%             \item formulate s.m.a.r.t.~goals
%             \item \textbf{how} we can do it
%             \item \textbf{when} will we do it
%         \end{itemize}
%     \end{itemize}
% \end{frame}



\begin{frame}{The Double Perspective of the Project Manager}

    \showmytikz{../shared/tikz/tikz-doubleperspective}

\end{frame}



% \begin{frame}[allowframebreaks]{Stake holders and Management}
  
%       \textbf{Someone} wants \textbf{something} to happen 
    
%     \begin{itemize}
%         \item Understand who wants the project to happen
%         \item What does this someone give you?
%         \item What does this someone expect?
%         \item What is this something?
%         \item How can it be achieved?
%         \item What do we know, and what do we not know?
%         \item How do we communicate and, if required, correct course?
%         \item Who may intervene? How do you react?
%     \end{itemize}
  
%     \framebreak 
    
%     Complex projects have \textbf{many stake holders.} 
    
%     \begin{itemize}
%         \item inside vs. outside
%         \item powerful vs. not powerful
%     \end{itemize}
    
    
%     $\rightarrow$ `Stakeholder Matrix'
%     \framebreak
    
%     \begin{itemize}
%         \item `above', such as project sponsors, project steering team, top level management and authorities, customers...
%         \item `below', such as the managers of sub-teams and  executing bodies
%         \item `on the same level', such as supporting line functions within the company
%         \item `outside' such as press etc.
%     \end{itemize}


%     \begin{tikzpicture}
%         \draw[<->] (-1.5,0) -- (1.5,0);
%         \draw[<->] (0,-1.5) -- (0,1.5);
%         \node[circle,draw,fill=white] at (0,0) {PM};
%     \end{tikzpicture}

%     Communication both to identify plans and to inform stake holders is a large part of project management.
% \end{frame}

